\phantomsection\addcontentsline{toc}{section}{Smart Grids}
\ECchapter{24}{Smart Grids}{Managing tomorrow's electricity networks}{ECPurple}

\ECObjectives{
\begin{itemize}
\item Explain why generation and demand must remain balanced.
\item Describe the roles of storage, forecasting and demand response.
\item Calculate a simple energy balance and identify resilience constraints.
\end{itemize}}

\begin{multicols}{2}
\begin{engineeringnews}
Electricity networks are changing from one-way systems dominated by large power stations into
interactive networks containing rooftop solar panels, wind farms, batteries, electric vehicles and
flexible loads. Digital control is becoming essential to coordinate these distributed resources.
\end{engineeringnews}

\begin{ecarticle}
An electricity grid must balance production and consumption at every moment. If generation is lower
than demand, system frequency tends to fall; if it is higher, frequency tends to rise. Traditional grids
adjusted a limited number of controllable power stations. Smart grids coordinate many smaller and more
variable devices.

Smart meters measure energy flows at short intervals. Forecasting software estimates demand and
renewable generation. Automated substations can reconfigure part of the network after a fault, while
batteries and flexible loads respond rapidly to imbalances. Demand response shifts selected uses in
time: for example, water heating or electric-vehicle charging may be delayed when the grid is heavily
loaded.

Distributed generation creates bidirectional power flows. A neighbourhood may import electricity in
the evening but export solar power at midday. Voltage can rise near photovoltaic installations, and
protection systems must distinguish normal reverse flow from a dangerous fault. Engineers therefore
combine power electronics, communication systems and network models.

Storage provides several services on different timescales. Batteries can smooth short fluctuations,
support frequency and reduce peak demand. Pumped-storage hydroelectricity can store much larger amounts
of energy for several hours or days. Electric vehicles may provide controlled charging or energy back
to the grid, but battery degradation, charging efficiency and user mobility must be respected.

Digitalisation increases exposure to cyberattack and communication failure. Local controllers must
retain safe behaviour when the central platform is unavailable. Resilience also requires physical
redundancy, black-start capability and plans for extreme weather. A smart grid is not merely connected;
it must remain stable, secure and understandable during abnormal conditions.
\end{ecarticle}

\begin{tcolorbox}[title=\sffamily\bfseries Did You Know?,colback=yellow!10,colframe=ECOrange]
Electricity travels through the network almost instantly, but energy planning covers seconds, hours,
seasons and years. No single storage technology is optimal for every timescale.
\end{tcolorbox}

\begin{engineeringfocus}
\textbf{Bidirectional smart grid}
\begin{center}
\resizebox{0.96\linewidth}{!}{%
\begin{tikzpicture}[font=\scriptsize,>=Latex,node distance=4mm]
\node[draw=ECPurple,fill=ECPurple!10,rounded corners] (grid) {distribution grid};
\node[draw=ECOrange,fill=ECOrange!10,rounded corners,left=of grid] (solar) {solar and wind};
\node[draw=ECGreen,fill=ECGreen!10,rounded corners,right=of grid] (loads) {homes and industry};
\node[draw=ECBlue,fill=ECLightBlue,rounded corners,below=of grid] (store) {storage and EVs};
\node[draw=ECRed,fill=ECRed!8,rounded corners,above=of grid] (control) {forecast and control};
\draw[<->,thick] (solar)--(grid); \draw[<->,thick] (grid)--(loads); \draw[<->,thick] (grid)--(store);
\draw[<->,thick,dashed] (control)--(grid);
\end{tikzpicture}}
\end{center}
\end{engineeringfocus}

\begin{keyfigures}
\begin{tabularx}{\linewidth}{@{}Xr@{}}
Instant requirement & power balance\\
Fast flexibility & batteries\\
Long-duration storage & pumped hydro\\
Flexible demand & smart charging\\
Critical risk & cyberattack
\end{tabularx}
\end{keyfigures}

\begin{tcolorbox}[title=\sffamily\bfseries Daily balancing example,colback=white,colframe=ECPurple]
\begin{center}
\begin{tikzpicture}
\begin{axis}[width=0.96\linewidth,height=48mm,xlabel={Hour},ylabel={Power (MW)},grid=major,
legend style={font=\scriptsize,at={(0.5,-0.22)},anchor=north,legend columns=3},ticklabel style={font=\scriptsize}]
\addplot[thick] table[x=hour,y=demand,col sep=comma]{data/EC24/load.csv};
\addplot[thick] table[x=hour,y=solar,col sep=comma]{data/EC24/load.csv};
\addplot[thick,dashed] table[x=hour,y=battery,col sep=comma]{data/EC24/load.csv};
\legend{Demand,Solar,Battery output}
\end{axis}
\end{tikzpicture}
\end{center}
\end{tcolorbox}

\begin{tcolorbox}[title=\sffamily\bfseries Quantitative application,colback=white,colframe=ECBlue]
Between 18:00 and 20:00, a school consumes a constant power of $140\,\mathrm{kW}$. Its photovoltaic
system produces $35\,\mathrm{kW}$, and a battery initially contains $180\,\mathrm{kWh}$ of usable
energy. The round-trip efficiency is not considered in this first estimate.
\begin{enumerate}
\item Calculate the power that must be supplied by the battery if no power is imported from the grid.
\item Determine the energy required during the two-hour period.
\item Decide whether the battery can cover the complete period.
\item Repeat the conclusion with a discharge efficiency of 90\% and a minimum reserve of $30\,\mathrm{kWh}$.
\end{enumerate}
\end{tcolorbox}

\begin{vocabulary}
\begin{tabularx}{\linewidth}{@{}lX@{}}
smart meter & compteur communicant\\
demand response & effacement / flexibilité\\
distributed generation & production décentralisée\\
peak demand & pointe de consommation\\
frequency control & réglage de fréquence\\
bidirectional flow & flux bidirectionnel\\
substation & poste électrique\\
black start & redémarrage autonome\\
load shedding & délestage\\
grid resilience & résilience du réseau
\end{tabularx}
\end{vocabulary}

\begin{questions}
\item Why must generation equal demand continuously?
\item How does demand response help the network?
\item What problem can midday solar export create?
\item Compare batteries and pumped-storage hydroelectricity.
\item Why should local controllers work without the central platform?
\end{questions}

\begin{engineeringchallenge}
Plan the energy management of a school with rooftop solar panels, a $200\,\mathrm{kWh}$ battery and ten
EV chargers. Define priorities for a normal day, a winter peak and a two-hour grid outage. Include one
rule that protects battery lifetime.
\end{engineeringchallenge}

\begin{speaking}
Should electricity suppliers be allowed to delay domestic appliances automatically during peak demand?
Discuss consent, comfort, price, privacy and grid stability.
\end{speaking}
\end{multicols}

\subsection*{Sources}
\ECsource{International Energy Agency -- Smart Grids}{https://www.iea.org/energy-system/electricity/smart-grids}
\ECsource{ENTSO-E}{https://www.entsoe.eu/}